% ---------------------------------------------------------------
% TEMPLATE TCC1 - SENAC SANTO AMARO
% ---------------------------------------------------------------
% Trabalho de Conclusão de Curso 1
% Elaborado para apresentação da proposta inicial do TCC
% Estrutura: 4 capítulos (Introdução, Revisão Bibliográfica,
%            Desenvolvimento, Referências Bibliográficas)
% ---------------------------------------------------------------

\documentclass[
    12pt,                % tamanho da fonte
    openright,           % capítulos começam em página ímpar
    oneside,             % para impressão em apenas um lado do papel
    a4paper,             % tamanho do papel
    english,             % idioma adicional
    brazil               % idioma principal
]{abntex2}

% ---------------------------------------------------------------
% PACOTES
% ---------------------------------------------------------------
\usepackage{lmodern}                % Fonte Latin Modern
\usepackage[T1]{fontenc}            % Seleção de códigos de fonte
\usepackage[utf8]{inputenc}         % Codificação do documento
\usepackage{lastpage}               % Usado pela Ficha catalográfica
\usepackage{indentfirst}            % Indenta o primeiro parágrafo
\usepackage{color}                  % Controle de cores
\usepackage{graphicx}               % Inclusão de gráficos
\usepackage{microtype}              % Melhorias de justificação
\usepackage{float}                  % Para posicionamento de figuras
\usepackage{booktabs}               % Tabelas profissionais
\usepackage{multirow}               % Células mescladas em tabelas
\usepackage{longtable}              % Tabelas longas
\usepackage{amsmath,amssymb}        % Símbolos matemáticos
\usepackage{pgfgantt}               % Diagramas de Gantt
\usepackage[brazilian,hyperpageref]{backref}  % Páginas com citações
\usepackage[alf]{abntex2cite}       % Citações ABNT

% ---------------------------------------------------------------
% CONFIGURAÇÕES DO PDF
% ---------------------------------------------------------------
\makeatletter
\hypersetup{
    pdftitle={\@title},
    pdfauthor={\@author},
    pdfsubject={TCC1 - SENAC Santo Amaro},
    pdfcreator={LaTeX with abnTeX2},
    pdfkeywords={tcc}{senac}{santo amaro}{engenharia de dados}{ciclo de vida dos dados}{pipeline de dados},
    colorlinks=true,
    linkcolor=blue,
    citecolor=blue,
    filecolor=magenta,
    urlcolor=blue,
    bookmarksdepth=4
}
\makeatother

% ---------------------------------------------------------------
% CONFIGURAÇÕES DE APARÊNCIA DO PDF FINAL
% ---------------------------------------------------------------
\definecolor{blue}{RGB}{41,5,195}

% ---------------------------------------------------------------
% INFORMAÇÕES DO DOCUMENTO
% ---------------------------------------------------------------
\titulo{Arquitetura para Gestão e Automação do Ciclo de Vida dos Dados em Ambientes Corporativos}
\autor{Giovanna Paiva Alves}
\local{São Paulo}
\data{2025}
\orientador{Prof. Nome do Orientador}

\instituicao{%
  SENAC -- Serviço Nacional de Aprendizagem Comercial
  \par
  Unidade Santo Amaro
  \par
  Curso de Tecnologia em Sistemas para Internet
}

\tipotrabalho{Trabalho de Conclusão de Curso I}

\preambulo{Proposta de Trabalho de Conclusão de Curso apresentado ao SENAC Santo Amaro como requisito parcial para aprovação na disciplina de TCC1 do curso de Tecnologia em Sistemas para Internet.}

% ---------------------------------------------------------------
% COMPILA O ÍNDICE
% ---------------------------------------------------------------
\makeindex

% ---------------------------------------------------------------
% INÍCIO DO DOCUMENTO
% ---------------------------------------------------------------
\begin{document}

\selectlanguage{brazil}
\frenchspacing

% ---------------------------------------------------------------
% ELEMENTOS PRÉ-TEXTUAIS
% ---------------------------------------------------------------
\pretextual
\imprimircapa
\imprimirfolhaderosto

% ---------------------------------------------------------------
% RESUMO
% ---------------------------------------------------------------
\begin{resumo}
Este documento apresenta a proposta inicial do Trabalho de Conclusão de Curso (TCC1) desenvolvido no SENAC Santo Amaro. O trabalho tem como objetivo desenvolver uma arquitetura para gestão e automação do ciclo de vida dos dados em ambientes corporativos, abrangendo as etapas de ingestão, transformação, armazenamento, versionamento, monitoramento e disponibilização para consumo analítico. A justificativa baseia-se na ausência, em muitos cenários corporativos, de estruturas padronizadas que organizem adequadamente o fluxo de dados, gerando inconsistências, baixa rastreabilidade e dificuldades na utilização dos dados para análise e inteligência artificial. A revisão bibliográfica abrange conceitos de ciclo de vida dos dados, engenharia de dados, arquiteturas de pipelines e metodologias de pesquisa em Design Science Research (DSR). Como solução proposta, pretende-se definir uma arquitetura de referência e desenvolver um protótipo funcional, validado por meio de um estudo de caso com dados reais ou simulados. Este trabalho está organizado em três capítulos que apresentam a introdução ao tema, a revisão bibliográfica que fundamenta a pesquisa e o desenvolvimento com a solução proposta e o cronograma de trabalho, seguidos das referências bibliográficas.

\vspace{\onelineskip}

\noindent
\textbf{Palavras-chave}: ciclo de vida dos dados. engenharia de dados. pipeline de dados. automação. arquitetura de dados. ambientes corporativos.
\end{resumo}

% ---------------------------------------------------------------
% SUMÁRIO
% ---------------------------------------------------------------
\pdfbookmark[0]{\contentsname}{toc}
\tableofcontents*
\cleardoublepage

% ---------------------------------------------------------------
% ELEMENTOS TEXTUAIS
% ---------------------------------------------------------------
\textual

% ===============================================================
% CAPÍTULO 1 - INTRODUÇÃO
% ===============================================================
\chapter{Introdução}
\label{cap:introducao}

A produção de dados nas organizações modernas cresce em volume, velocidade e variedade de forma acelerada. Sistemas transacionais, APIs, logs de aplicações e bancos de dados geram continuamente grandes quantidades de informações que precisam ser coletadas, processadas e disponibilizadas de forma confiável para apoiar decisões estratégicas. Contudo, a ausência de uma arquitetura bem definida para o gerenciamento dessas informações compromete a qualidade, a rastreabilidade e a utilidade dos dados ao longo de todo o seu ciclo de vida. Este trabalho propõe o desenvolvimento de uma arquitetura para gestão e automação do ciclo de vida dos dados em ambientes corporativos, contemplando desde a ingestão até o consumo analítico.

% ---------------------------------------------------------------
\section{Contexto}
\label{sec:contexto}

Organizações modernas produzem grandes volumes de dados provenientes de diversas fontes, como sistemas transacionais, APIs, logs de aplicações e bancos de dados. Em muitos cenários corporativos, esses dados são coletados, armazenados e processados sem uma estrutura padronizada que organize adequadamente seu ciclo de vida. Essa ausência de padronização gera consequências relevantes: inconsistências entre fontes de informação, retrabalho no processamento de dados, baixa rastreabilidade das transformações realizadas e dificuldades na utilização dos dados para análises e treinamento de modelos de inteligência artificial \cite{Shah2021}.

A gestão do ciclo de vida dos dados abrange as etapas de criação, armazenamento, uso, compartilhamento, arquivamento e eliminação das informações \cite{Shah2021}. Em ambientes corporativos, a complexidade cresce à medida que múltiplas fontes e sistemas distintos precisam ser integrados em fluxos coerentes de processamento. Além disso, regulamentações de proteção de dados, como a Lei Geral de Proteção de Dados (LGPD) no Brasil e o GDPR na Europa, impõem requisitos adicionais de controle, retenção e descarte seguro das informações \cite{Scope2022}.

Diante desse cenário, a engenharia de dados surge como área fundamental para o desenvolvimento de arquiteturas que organizem, automatizem e monitorem o fluxo de dados corporativos. A definição de pipelines de dados estruturados, com etapas bem delimitadas de ingestão, transformação, armazenamento e disponibilização, constitui uma resposta direta aos desafios identificados nas organizações contemporâneas.

% ---------------------------------------------------------------
\section{Justificativa}
\label{sec:justificativa}

A ausência de uma arquitetura padronizada para gestão do ciclo de vida dos dados representa um problema recorrente em organizações que lidam com múltiplas fontes de informação. Esse problema se manifesta em inconsistências entre sistemas, perda de rastreabilidade das transformações aplicadas aos dados e dificuldades na manutenção dos pipelines de processamento. Tais questões comprometem não apenas a qualidade das análises realizadas, mas também a capacidade das organizações de utilizar os dados para o desenvolvimento de modelos de inteligência artificial e de apoio à tomada de decisão.

Do ponto de vista acadêmico, a área de engenharia de dados ainda carece de frameworks de referência amplamente validados que abranjam de forma integrada todas as etapas do ciclo de vida dos dados em contextos corporativos. Embora existam trabalhos que abordam aspectos isolados, como limpeza de dados \cite{Valova2025}, gestão de conformidade \cite{Scope2022} e modelos de ciclo de vida para governos \cite{Shah2021}, uma arquitetura de referência voltada à automação e ao monitoramento de ponta a ponta em ambientes corporativos privados ainda representa uma lacuna relevante na literatura.

Do ponto de vista prático, a proposta deste trabalho visa contribuir com organizações que buscam estruturar seus fluxos de dados de forma mais eficiente, reduzindo retrabalho, aumentando a confiabilidade das informações e facilitando a manutenção dos pipelines ao longo do tempo.

% ---------------------------------------------------------------
\section{Objetivo}
\label{sec:objetivo}

\subsection{Objetivo principal}

Desenvolver uma arquitetura para gestão e automação do ciclo de vida dos dados, capaz de estruturar e padronizar os processos de ingestão, processamento, armazenamento e disponibilização de dados em ambientes corporativos.

\subsection{Objetivos secundários}

\begin{enumerate}
    \item Analisar os principais conceitos relacionados ao ciclo de vida dos dados, engenharia de dados e arquiteturas de pipelines de dados em ambientes corporativos;
    \item Identificar os principais desafios enfrentados por organizações na gestão e processamento de grandes volumes de dados provenientes de múltiplas fontes;
    \item Definir uma arquitetura de referência para organização do ciclo de vida dos dados, contemplando etapas como ingestão, transformação, armazenamento e disponibilização;
    \item Desenvolver um protótipo baseado na arquitetura proposta, capaz de automatizar pipelines de dados e apoiar a gestão das diferentes etapas do ciclo de vida dos dados;
    \item Aplicar a arquitetura em um estudo de caso, utilizando um conjunto de dados real ou simulado para demonstrar sua aplicabilidade e avaliar aspectos como automação do fluxo, rastreabilidade das transformações e facilidade de manutenção.
\end{enumerate}

% ===============================================================
% CAPÍTULO 2 - REVISÃO BIBLIOGRÁFICA
% ===============================================================
\chapter{Revisão Bibliográfica}
\label{cap:revisao}

A revisão bibliográfica tem como finalidade apresentar o embasamento teórico do trabalho, demonstrando o conhecimento da autora sobre o tema e identificando o estado da arte. Este capítulo está organizado em quatro seções: referencial teórico, metodologia da pesquisa bibliográfica, fundamentos técnicos e trabalhos relacionados.

% ---------------------------------------------------------------
\section{Referencial teórico}
\label{sec:referencial}

\subsection{Ciclo de Vida dos Dados}

O ciclo de vida dos dados compreende o conjunto de fases pelas quais uma informação passa, desde sua criação ou coleta até o seu descarte ou arquivamento definitivo. Em organizações orientadas a dados, esse ciclo envolve etapas como criação, armazenamento, processamento, compartilhamento, uso analítico, arquivamento e eliminação. \citeonline{Shah2021} identificaram, por meio de uma revisão sistemática da literatura, 76 modelos de ciclo de vida de dados, propondo um framework integrado para governos orientados a dados. Os autores destacam que o gerenciamento eficaz do ciclo de vida dos dados é um desafio crescente para organizações que lidam com grandes volumes de informações provenientes de fontes diversas.

Em ambientes corporativos, a gestão do ciclo de vida dos dados apresenta camadas adicionais de complexidade. A necessidade de integrar sistemas heterogêneos, manter a conformidade com regulamentações de proteção de dados e garantir a rastreabilidade das transformações aplicadas às informações exige a adoção de arquiteturas estruturadas e de mecanismos automatizados de controle \cite{Scope2022}.

\subsection{Engenharia de Dados e Pipelines}

A engenharia de dados é a disciplina responsável pelo desenvolvimento e manutenção de infraestruturas que permitem a coleta, transformação e disponibilização de dados de forma confiável e escalável. No centro dessa disciplina estão os pipelines de dados — fluxos automatizados que conectam fontes de dados a destinos analíticos, realizando ao longo do caminho transformações necessárias para garantir a qualidade e a utilidade das informações.

Arquiteturas de pipelines bem definidas possibilitam a padronização dos processos de processamento de dados, a automação de tarefas repetitivas e o monitoramento contínuo da qualidade e integridade das informações. A qualidade dos dados de entrada em um pipeline é fator determinante para a confiabilidade das análises produzidas, o que torna a etapa de limpeza e transformação crítica para o sucesso de qualquer arquitetura de dados \cite{Valova2025}.

\subsection{Governança e Conformidade de Dados}

A governança de dados envolve o conjunto de políticas, processos e tecnologias utilizados para garantir que os dados de uma organização sejam gerenciados de forma consistente, segura e em conformidade com as regulamentações aplicáveis. Em contextos corporativos, regulamentações como a LGPD e o GDPR impõem requisitos específicos de retenção e descarte de dados que precisam ser incorporados às arquiteturas de gerenciamento de informações.

\citeonline{Scope2022} propõem uma abordagem abrangente para integração de políticas de retenção e purga de dados em sistemas de gerenciamento de banco de dados, destacando que a solução deve ser transparente para o usuário final, exigir mínima configuração pelo administrador e impor overhead de desempenho negligenciável.

% ---------------------------------------------------------------
\section{Metodologia da pesquisa bibliográfica}
\label{sec:metodologia_pesquisa}

A pesquisa bibliográfica deste trabalho foi conduzida por meio do método do funil de pesquisa (\textit{Research Funnel}), uma abordagem sistemática que parte de buscas amplas e progressivamente aplica critérios de filtragem para selecionar os trabalhos mais relevantes ao tema investigado.

As buscas foram realizadas nas bases de dados Google Scholar, IEEE Xplore, ACM Digital Library e ResearchRabbit, utilizando as seguintes palavras-chave em português e inglês: \textit{data lifecycle management}, \textit{data pipeline}, \textit{data engineering}, \textit{data governance}, \textit{ETL automation}, \textit{ciclo de vida dos dados}, \textit{pipeline de dados} e \textit{engenharia de dados}.

Os critérios de inclusão adotados foram: publicações entre 2018 e 2025, nos idiomas português e inglês, que abordassem diretamente conceitos de ciclo de vida dos dados, pipelines de dados ou arquiteturas de gerenciamento de dados em contextos corporativos ou governamentais. Foram excluídos trabalhos que tratassem exclusivamente de domínios muito específicos sem relação com ambientes corporativos de tecnologia da informação, além de publicações sem revisão por pares.

O processo de seleção seguiu as etapas do funil: (1) identificação de publicações por palavras-chave; (2) triagem por título e resumo; (3) leitura completa dos trabalhos pré-selecionados; e (4) seleção final com base na relevância direta para o problema de pesquisa.

% ---------------------------------------------------------------
\section{Fundamentos}
\label{sec:fundamentos}

\subsection{Arquiteturas de Dados: Data Lake, Data Warehouse e Lakehouse}

As principais arquiteturas de armazenamento e processamento de dados utilizadas em ambientes corporativos incluem o \textit{Data Warehouse}, o \textit{Data Lake} e, mais recentemente, o \textit{Data Lakehouse}. O \textit{Data Warehouse} é uma solução estruturada e otimizada para consultas analíticas sobre dados históricos, enquanto o \textit{Data Lake} permite o armazenamento de dados brutos em seu formato original, com maior flexibilidade e menor custo. O \textit{Data Lakehouse} combina características das duas abordagens, unindo a flexibilidade do \textit{Data Lake} com as capacidades transacionais e de governança do \textit{Data Warehouse}.

\subsection{Orquestração de Pipelines}

A orquestração de pipelines de dados refere-se ao gerenciamento automatizado do sequenciamento, agendamento e monitoramento das tarefas de processamento de dados. Ferramentas como Apache Airflow, Prefect e Dagster são amplamente utilizadas para esse fim, permitindo a definição de fluxos de trabalho como grafos acíclicos dirigidos (DAGs), nos quais cada nó representa uma tarefa e as arestas representam dependências entre elas.

\subsection{Design Science Research (DSR) como Metodologia de Desenvolvimento}

A metodologia Design Science Research (DSR) orienta pesquisas voltadas à criação de artefatos — como sistemas, modelos, frameworks e algoritmos — que solucionem problemas identificados em contextos reais. A DSR é particularmente adequada a trabalhos de engenharia de software e sistemas de informação, pois propõe um ciclo iterativo de design, construção e avaliação do artefato produzido. Neste trabalho, a DSR orientará o desenvolvimento da arquitetura proposta e do protótipo a ser implementado, com avaliação baseada em estudo de caso.

% ---------------------------------------------------------------
\section{Trabalhos relacionados}
\label{sec:trabalhos_relacionados}

Esta seção apresenta três trabalhos que compartilham o foco central desta pesquisa: a gestão estruturada do ciclo de vida dos dados em contextos organizacionais. A tabela comparativa ao final sintetiza as principais diferenças entre as abordagens.

\subsection{Trabalho 1: DaLiF -- A Data Lifecycle Framework for Data-Driven Governments}

\citeonline{Shah2021} abordaram o problema da ausência de um framework consolidado para gestão do ciclo de vida de dados em governos orientados a dados. Para resolvê-lo, os autores realizaram uma Revisão Sistemática da Literatura (RSL) de 76 modelos de ciclo de vida de dados e propuseram o framework DaLiF, voltado especificamente ao setor público. Os resultados demonstraram que o framework é capaz de cobrir as principais dimensões da gestão de dados governamentais. No entanto, a solução apresenta como limitação o foco restrito ao setor público, sem contemplar as particularidades e os requisitos tecnológicos dos ambientes corporativos privados. O presente trabalho diferencia-se ao propor uma arquitetura orientada à automação de pipelines em ambientes corporativos privados, com validação por meio de protótipo funcional.

\subsection{Trabalho 2: Compliance and Data Lifecycle Management in Databases and Backups}

\citeonline{Scope2022} abordaram o problema de gestão de conformidade legal no ciclo de vida dos dados armazenados em sistemas de banco de dados, especialmente no contexto de regulamentações como HIPAA e GDPR. Os autores propuseram a integração de políticas de retenção e purga diretamente nos sistemas de gerenciamento de banco de dados (SGBDs), com foco na transparência para o usuário e no mínimo impacto de desempenho. Os resultados indicaram viabilidade técnica da abordagem. A limitação principal está no escopo restrito à camada de banco de dados, sem contemplar as etapas de ingestão, transformação e disponibilização analítica dos dados. Este trabalho avança ao propor uma arquitetura que abrange todo o ciclo de vida dos dados, incorporando os requisitos de conformidade de forma transversal ao pipeline.

\subsection{Trabalho 3: Data Cleaning -- Methods, Challenges, and Practices for Improving Data Quality}

\citeonline{Valova2025} abordaram o problema da qualidade dos dados na etapa de pré-processamento de pipelines analíticos, avaliando ferramentas e técnicas de limpeza de dados aplicadas a fontes heterogêneas, incluindo dispositivos IoT e datasets estruturados. Os autores propuseram um modelo conceitual de pipeline de limpeza de dados, avaliando ferramentas como Pandas, OpenRefine e Trifacta. Os resultados mostraram que cada ferramenta apresenta forças e limitações específicas a depender do volume e da natureza dos dados. A principal limitação é o foco isolado na etapa de limpeza, sem integração com as demais fases do ciclo de vida. O presente trabalho diferencia-se ao tratar a qualidade dos dados como parte integrante de uma arquitetura mais ampla, que automatiza e monitora todo o fluxo de processamento.

\vspace{0.5cm}

\begin{table}[htb]
    \centering
    \caption{Comparação entre trabalhos relacionados}
    \label{tab:trabalhos_relacionados}
    \begin{tabular}{l|c|c|c|c}
        \toprule
        \textbf{Critério} & \textbf{Shah (2021)} & \textbf{Scope (2022)} & \textbf{Valova (2025)} & \textbf{Este trabalho} \\
        \midrule
        Foco              & Ciclo de vida & Conformidade  & Qualidade/Limpeza & Ciclo de vida completo \\
        \midrule
        Domínio           & Setor público & BD/SGBD       & IoT/Analytics     & Corporativo privado \\
        \midrule
        Metodologia       & RSL           & Proposta técn.& Experimento       & DSR + Estudo de caso \\
        \midrule
        Protótipo         & Não           & Parcial       & Sim               & Sim \\
        \midrule
        Automação         & Não           & Parcial       & Não               & Sim \\
        \bottomrule
    \end{tabular}
    \legend{Fonte: Elaborado pela autora (2025)}
\end{table}

% ===============================================================
% CAPÍTULO 3 - DESENVOLVIMENTO
% ===============================================================
\chapter{Desenvolvimento}
\label{cap:desenvolvimento}

Neste capítulo, apresenta-se a solução proposta para o problema identificado, bem como o cronograma de trabalho previsto para o desenvolvimento completo do projeto no TCC2.

% ---------------------------------------------------------------
\section{Solução proposta}
\label{sec:solucao}

A solução proposta consiste em uma arquitetura de referência para gestão e automação do ciclo de vida dos dados em ambientes corporativos, acompanhada de um protótipo funcional que valide os conceitos definidos. A arquitetura será desenvolvida seguindo os princípios da metodologia Design Science Research (DSR), com ciclos iterativos de design, implementação e avaliação.

\begin{itemize}
    \item \textbf{Visão geral da solução}: será desenvolvida uma arquitetura de pipelines de dados que organize o fluxo de processamento desde a origem (sistemas transacionais, APIs, arquivos) até o consumo final (camadas analíticas, dashboards, modelos de IA), com mecanismos de rastreabilidade e monitoramento em cada etapa;

    \item \textbf{Arquitetura}: a solução contemplará as camadas de ingestão (\textit{raw/landing zone}), processamento e transformação (\textit{staging/trusted zone}) e disponibilização analítica (\textit{refined/delivery zone}), seguindo o padrão de arquitetura medalhão (\textit{Bronze, Silver, Gold}). A orquestração dos fluxos será realizada por meio de uma ferramenta dedicada, como Apache Airflow ou Prefect;

    \item \textbf{Tecnologias e ferramentas}: serão utilizadas tecnologias amplamente empregadas na área de engenharia de dados, como Python para transformações, Apache Spark ou dbt para processamento, armazenamento em camadas baseado em data lake (ex.: MinIO ou AWS S3), e ferramentas de orquestração como Apache Airflow. A escolha final será definida com base nos requisitos do estudo de caso;

    \item \textbf{Funcionalidades principais}: ingestão automatizada de dados de múltiplas fontes; transformação e padronização seguindo regras definidas; versionamento dos dados e das transformações aplicadas; monitoramento de qualidade dos dados em cada camada; registro de linhagem dos dados (\textit{data lineage}); e disponibilização para consumo analítico;

    \item \textbf{Metodologia de desenvolvimento}: a abordagem adotada será a Design Science Research (DSR), com fases de identificação do problema, definição dos objetivos, design e desenvolvimento do artefato, demonstração por meio do protótipo e avaliação dos resultados;

    \item \textbf{Critérios de validação}: a solução será avaliada por meio de um estudo de caso com dados reais ou simulados, analisando métricas como grau de automação do fluxo de dados, cobertura das etapas do ciclo de vida, rastreabilidade das transformações, facilidade de manutenção do pipeline e qualidade dos dados ao longo das camadas.
\end{itemize}

% ---------------------------------------------------------------
\section{Cronograma de trabalho}
\label{sec:cronograma}

O cronograma a seguir apresenta, no formato de diagrama de Gantt, as atividades previstas para o desenvolvimento do TCC2, distribuídas ao longo das semanas do semestre.

\begin{figure}[htb]
    \centering
    \caption{Cronograma de atividades -- Diagrama de Gantt}
    \label{fig:gantt}
    \begin{ganttchart}[
        hgrid,
        vgrid,
        x unit=0.72cm,
        y unit title=0.6cm,
        y unit chart=0.7cm,
        title height=1,
        bar height=0.5,
        bar label font=\scriptsize,
        title label font=\small,
        milestone label font=\scriptsize\itshape,
        group label font=\scriptsize\bfseries,
        bar/.append style={fill=blue!40},
        group/.append style={draw=black, fill=blue!70},
        milestone/.append style={fill=red!70},
    ]{1}{16}
        \gantttitle{Semanas}{16} \\
        \gantttitlelist{1,...,16}{1} \\

        \ganttgroup{Fase 1: Planejamento}{1}{4} \\
        \ganttbar{Revisão bibliográfica complementar}{1}{3} \\
        \ganttbar{Definição da arquitetura de referência}{2}{4} \\
        \ganttbar{Definição do estudo de caso}{3}{4} \\
        \ganttmilestone{Arquitetura definida}{4} \\

        \ganttgroup{Fase 2: Desenvolvimento}{5}{11} \\
        \ganttbar{Implementação camada de ingestão}{5}{7} \\
        \ganttbar{Implementação camada de transformação}{7}{9} \\
        \ganttbar{Orquestração e monitoramento}{9}{11} \\
        \ganttmilestone{Protótipo funcional}{11} \\

        \ganttgroup{Fase 3: Validação e Escrita}{12}{16} \\
        \ganttbar{Aplicação do estudo de caso}{12}{14} \\
        \ganttbar{Avaliação de métricas}{13}{15} \\
        \ganttbar{Redação final do TCC2}{13}{16} \\
        \ganttmilestone{Entrega final}{16}

        \ganttlink{elem2}{elem3}
        \ganttlink{elem3}{elem4}
        \ganttlink{elem4}{elem5}
        \ganttlink{elem7}{elem8}
        \ganttlink{elem8}{elem9}
        \ganttlink{elem9}{elem10}
        \ganttlink{elem12}{elem13}
    \end{ganttchart}
    \legend{Fonte: Elaborado pela autora (2025)}
\end{figure}

% ---------------------------------------------------------------
% ELEMENTOS PÓS-TEXTUAIS
% ---------------------------------------------------------------
\postextual

% ===============================================================
% REFERÊNCIAS BIBLIOGRÁFICAS
% ===============================================================
\bibliography{bibliografia}

\end{document}
